\documentclass[12pt]{article}
\usepackage[utf8]{inputenc}
\usepackage[T1]{fontenc}
\usepackage[margin=1in]{geometry}
\usepackage{setspace}
\usepackage{xcolor}
\usepackage[colorlinks=true, urlcolor=blue!60!black, linkcolor=blue!60!black]{hyperref}
\setlength{\parindent}{2em}
\setlength{\parskip}{1.25em}
\renewcommand{\baselinestretch}{1.0}
\begin{document}
\begin{center}
{\LARGE \textit{United States Steel Corp. v. Multistate Tax Commission}, 434 U.S. 452 (1978)\par}
\vspace{0.8em}
{\large Case Brief\par}
\end{center}
\vspace{1em}

\section{Memory Jogger}
Applying the functional test of \textit{Virginia v. Tennessee}, the Court held 7--2 that the Multistate Tax Compact needed no congressional consent because it gave member states no power they lacked individually and did not enhance state power at the expense of federal supremacy (\href{https://supreme.justia.com/cases/federal/us/434/452/}{Justia}).

\section{Detailed Case Facts}
After the Supreme Court held in \textit{Northwestern States Portland Cement Co. v. Minnesota}, 358 U.S. 450 (1959), that states may tax fairly apportioned net income from a foreign corporation's interstate operations, Congress enacted minimum nexus rules in Pub. L. No. 86-272 and commissioned a federal study (the Willis Committee) to recommend uniform state tax legislation (\href{https://www.law.cornell.edu/supremecourt/text/434/452}{Cornell LII Opinion}). When federal bills stalled, state tax administrators drafted the Multistate Tax Compact in 1966 to facilitate multistate tax administration, promote uniformity, improve taxpayer compliance, and avoid duplicative taxation. The Compact took effect on August 4, 1967, upon adoption by seven states; 21 states had joined by 1972 (and 19 were members at decision after four withdrew and two joined). Twelve bills seeking congressional consent were introduced in Congress, but none passed.

Article VI of the Compact created the Multistate Tax Commission, composed of member-state tax administrators, to study state tax systems and recommend uniform rules. Under Article VII, the Commission's uniform regulations were strictly advisory; each member state retained complete authority to reject, amend, or disregard them. Under optional Article VIII, a member state could request that the Commission conduct interstate corporate audits on its behalf and seek compulsory judicial process in any adopting state's courts, while each individual state retained sole power over tax rates, tax bases, assessments, and collection. Article X permitted any member state to withdraw unilaterally by enacting a repealing statute (\href{https://supreme.justia.com/cases/federal/us/434/452/}{Justia Opinion}).

\section{Procedural History}
In 1972, United States Steel Corp. and other major corporations filed a class action in the Southern District of New York on behalf of multistate taxpayers threatened with Commission audits, naming the Commission, its commissioners, and its executive director as defendants. Plaintiffs sought declaratory and injunctive relief, alleging that the Compact violated the Compact Clause for lack of congressional consent, burdened interstate commerce, and infringed the Fourth and Fourteenth Amendments. After denying a motion to dismiss (367 F. Supp. 107 (S.D.N.Y. 1973)), a three-judge District Court granted summary judgment for defendants (417 F. Supp. 795 (S.D.N.Y. 1976)). Abandoning their Fourth Amendment claim, plaintiffs took a direct appeal to the Supreme Court.

\section{Judicial Votes}
The Supreme Court \textbf{affirmed} the three-judge District Court by a \textbf{7--2} vote. The case was argued on October 11, 1977, and decided on February 21, 1978 (\href{https://supreme.justia.com/cases/federal/us/434/452/}{Justia}; \href{https://www.law.cornell.edu/supremecourt/text/434/452}{Cornell LII}):

\begin{itemize}
    \item \textbf{Majority Opinion:} Justice Powell delivered the opinion of the Court, joined by Chief Justice Burger and Justices Brennan, Stewart, Marshall, Rehnquist, and Stevens.
    \item \textbf{Concurring Opinions:} None.
    \item \textbf{Dissenting Opinion:} Justice White filed a dissenting opinion, joined by Justice Blackmun.
\end{itemize}

\section{Holding}
\textbf{Question Presented:}

Whether the Multistate Tax Compact---a multilateral interstate agreement creating an ongoing administrative commission without congressional approval---is invalid under the Compact Clause, U.S. Const. art. I, \S\ 10, cl. 3, or under the Commerce Clause and the Fourteenth Amendment.

\textbf{Holding:} No. In a 7--2 decision, the Supreme Court affirmed summary judgment and held that the Multistate Tax Compact does not require congressional consent under the Compact Clause (\href{https://supreme.justia.com/cases/federal/us/434/452/}{Justia Opinion}). Reaffirming the functional rule of \href{https://supreme.justia.com/cases/federal/us/148/503/}{\textit{Virginia v. Tennessee}}, 148 U.S. 503 (1893), the Court held that congressional approval is required only for interstate agreements directed to forming a combination tending to increase state political power in a way that may encroach upon or interfere with the just supremacy of the United States. The Court also rejected the Commerce Clause and Fourteenth Amendment challenges.

\section{Analysis of Opinions}
\subsection{Majority Opinion (Justice Powell)}

\subsection{Rejecting a Literal Reading of the Compact Clause}
Justice Powell began by observing that a literal reading of Article I, \S\ 10, cl. 3 would require congressional approval for every trivial agreement between states, regardless of subject or federal interest (\href{https://www.law.cornell.edu/supremecourt/text/434/452}{Cornell LII Opinion}). Surveying the Articles of Confederation, the Constitutional Convention, Justice Story's \textit{Commentaries}, and early decisions such as \textit{Holmes v. Jennison}, 14 Pet. 540 (1840), Powell concluded that the Framers' precise distinctions among ``treaty,'' ``compact,'' and ``agreement'' had long been lost. Instead, constitutional doctrine settled on Justice Field's functional test in \href{https://supreme.justia.com/cases/federal/us/148/503/}{\textit{Virginia v. Tennessee}}, 148 U.S. at 519, applied as a holding in \textit{New Hampshire v. Maine}, 426 U.S. 363 (1976): congressional consent is required only when an agreement tends to increase state political power so as to encroach upon federal supremacy.

Powell rejected appellants' effort to confine \textit{Virginia v. Tennessee} to bilateral agreements lacking an administrative agency. While agreeing that the proper inquiry focuses on \textit{potential} rather than actual impact on federal supremacy---and applies equally to formal compacts and informal reciprocal legislation---the Court held that neither multilateral membership nor an ongoing administrative body triggers the Compact Clause by itself. What matters is the substantive impact on the federal structure.

\subsection{Applying the Functional Test to the Multistate Tax Compact}
Measured against federal supremacy, the Compact failed to trigger the consent requirement for three structural reasons:
\begin{itemize}
  \item \textbf{No Enhancement of State Power \textit{Quoad} the National Government:} Even if collective action increased states' bargaining leverage against multistate corporations, the test is whether the Compact enhances state power relative to the Federal Government. The Compact authorized no power that member states could not exercise individually, delegated no sovereign authority to the Commission, left every state free to adopt or reject Commission rules, and allowed unilateral withdrawal at any time.
  \item \textbf{No Encroachment on Interstate Commerce or Foreign Relations:} Appellants pointed to unitary-business audits, broad definitions of apportionable business income, affiliate-data filings, reciprocal court-process enforcement under Article VIII, and multinational audits. Because each member state could constitutionally adopt those same taxing methods and reciprocal subpoena laws on its own, coordinated administration did not aggrandize state authority at federal expense; Congress retained plenary power to preempt the field or ratify a contrary tax treaty.
  \item \textbf{No Impairment of Nonmember States' Sovereignty:} Any economic pressure on nonmember states to conform their tax rules stemmed from lawful fiscal competition among states, not from an unconstitutional compact.
\end{itemize}
In Footnote 33, Powell answered the dissent directly: the mere existence of a ``federal interest'' in interstate taxation or enhanced state lobbying clout in Congress does not threaten federal supremacy where Congress's legislative authority remains plenary (\href{https://www.law.cornell.edu/supremecourt/text/434/452}{Cornell LII Opinion}). Finally, the Court dismissed appellants' Commerce Clause and Fourteenth Amendment harassment allegations as unpreserved, unsupported by the record, and irrelevant to the Compact's facial validity because only individual states---not the Commission---could issue tax assessments.

\subsection{Dissenting Opinion (Justice White)}

Justice White accused the majority of paying lip service to ``potential'' encroachment while demanding proof of actual harm, stripping the Compact Clause of independent force beyond the Dormant Commerce Clause (\href{https://supreme.justia.com/cases/federal/us/434/452/}{Justia Opinion}). Emphasizing that Congress had refused to enact twelve compact-consent bills and that the Commission openly lobbied to ``bottle up'' federal Willis Bills and oppose a pending U.S.--U.K. tax treaty, White argued that both Congress and the President had unmistakably demonstrated a federal sovereign interest. He also highlighted synergistic powers beyond isolated state action---Commission-initiated audits, mandatory apportionment options, an all-or-nothing reciprocal subpoena bloc, and a seven-state activation trigger---and stressed that the Compact Clause protects nonmember states from competitive coercion as well as federal supremacy. Calling the Clause ``conciliatory rather than prohibitive,'' White warned that upholding so pervasive a multistate institution without consent left ``very little life remaining'' in the Compact Clause.

\subsection{Subsequent Developments}

\begin{itemize}

\item \textbf{Applied to Regional Banking Statutes:} In \href{https://supreme.justia.com/cases/federal/us/472/159/}{\textit{Northeast Bancorp, Inc. v. Board of Governors}}, 472 U.S. 159 (1985), the Court applied \textit{U.S. Steel} to uphold regional New England banking statutes, refining the indicia of a ``compact'' (such as a joint organization, non-unilateral withdrawal, and reciprocal obligations) while confirming that even if an agreement existed, it neither enhanced state power over the Federal Government nor impaired nonmember states' sovereignty.

\end{itemize}

\section{Examples: Future Applications}
\begin{enumerate}

\item \textbf{Same-Side (No Consent Needed): An Advisory Interstate AI Oversight Commission.} Fifteen states enact an unconsented ``Interstate AI Oversight Compact'' creating a joint commission that drafts model algorithmic-transparency rules and conducts voluntary joint technical audits of AI companies at a member state's request. Member states may accept or reject any model rule, impose their own fines, and withdraw upon 30 days' notice. Under \textit{U.S. Steel}, the compact is valid without congressional consent. Even though Congress has debated federal AI legislation (a ``federal interest''), the compact delegates no sovereign authority, exercises no power unavailable to individual states, permits unilateral withdrawal, and leaves Congress's Commerce Clause preemption power intact.

\item \textbf{Same-Side (No Consent Needed): Reciprocal Recognition of Teaching Licenses.} Twenty States enact identical statutes agreeing to recognize one another's teaching licenses and creating an advisory council of state education officials that recommends common certification standards. Each State keeps full control over whom it licenses, may ignore the council's recommendations, and may withdraw by repealing its statute. A teachers' union argues that the arrangement is an unconsented compact. Under \textit{U.S. Steel}, the relevant question is the agreement's potential impact on federal supremacy, and here there is none: every State could recognize out-of-state licenses on its own, the council exercises no sovereign power, and Congress remains free to legislate. The arrangement therefore requires no congressional consent, whether it is called a compact or reciprocal legislation.

\item \textbf{Opposite-Side (Consent Required): A Binding Atlantic Maritime Security Compact.} Suppose instead twelve coastal states form an unconsented ``Atlantic Maritime Security Compact'' that establishes a binding regional port authority with independent power to levy harbor tariffs on foreign vessels, negotiate bilateral customs-inspection agreements with European nations, and prohibit any member state from withdrawing for 25 years without unanimous consent. Unlike the Multistate Tax Compact in \textit{U.S. Steel}, this pact violates the Compact Clause (and the Treaty and Import--Export/Tonnage Clauses). It binds member states irrevocably, delegates sovereign taxing and enforcement powers to an interstate body that no single state could lawfully exercise alone, and directly encroaches upon exclusive federal supremacy over foreign affairs and foreign commerce (\href{https://www.law.cornell.edu/constitution-conan/article-1/section-1/clause-3/requirement-of-congressional-consent-to-compacts}{Cornell LII Constitution Annotated}). Under \textit{Virginia v. Tennessee} and \textit{U.S. Steel}, such an agreement cannot operate without express congressional consent.

\item \textbf{Opposite-Side (Consent Required): Transferring a County Between States.} Two neighboring States agree that a populous county will be detached from the first State and annexed to the second, and each legislature approves the transfer. The second State would gain population, and with it potential seats in the House of Representatives and electoral votes. Under the functional test of \textit{Virginia v. Tennessee} that \textit{U.S. Steel} reaffirmed, an agreement requires consent when it tends to increase the political power of the States in a way that may encroach upon the just supremacy of the United States. Shifting a block of population and territory between States alters representation in the national government itself, so the agreement cannot take effect without Congress's consent.

\item \textbf{Fence-Sitter (Unsettled): The National Popular Vote Interstate Compact.} States agree to award their presidential electors to the winner of the national popular vote, effective only once States holding a majority of electoral votes have joined. Supporters rely on \textit{U.S. Steel}: each State already has plenary power under Article II to choose how to appoint its electors, the agreement creates no commission with sovereign power, and members may withdraw. Critics answer with Justice White's point that the Clause must forbid some joint actions that individual States could take alone, and argue that a bloc controlling the outcome of presidential elections changes the national political structure and dilutes the influence of nonmember States. Because \textit{U.S. Steel} measured harm only against federal supremacy and the Court has never applied it to a compact affecting the selection of federal officers, the outcome is genuinely uncertain.

\end{enumerate}

\section{Critique}
\textbf{Nullification of Constitutional Text:} Originalist and textualist critics echo Justice White: \textit{U.S. Steel} reads ``any Agreement or Compact'' so narrowly that virtually no domestic interstate arrangement ever requires consent unless it already violates another constitutional limit.

\textbf{Blindness to Horizontal Federalism:} By testing state power exclusively \textit{quoad} the Federal Government, Powell's opinion gives short shrift to the danger that a regional cartel of states can impose political or economic externalities on excluded sister states without directly usurping a federal statute.

\textbf{Underestimating Institutional Synergy:} Treating a permanent, staffed multistate agency as constitutionally identical to 21 states acting independently ignores collective-action dynamics and coordinated lobbying power.

\section{Key Quotations}
\begin{quote}
``Read literally, the Compact Clause would require the States to obtain congressional approval before entering into any agreement among themselves, irrespective of form, subject, duration, or interest to the United States.''
\end{quote}
\noindent --- Justice Powell, \href{https://supreme.justia.com/cases/federal/us/434/452/}{434 U.S. at 459}.

\begin{quote}
``The Clause reaches both `agreements' and `compacts,' the formal as well as the informal. The relevant inquiry must be one of impact on our federal structure.''
\end{quote}
\noindent --- Justice Powell, \href{https://supreme.justia.com/cases/federal/us/434/452/}{434 U.S. at 470--471}.

\begin{quote}
``But the test is whether the Compact enhances state power \textit{quoad} the National Government. This pact does not purport to authorize the member States to exercise any powers they could not exercise in its absence. Nor is there any delegation of sovereign power to the Commission; each State retains complete freedom to adopt or reject the rules and regulations of the Commission.''
\end{quote}
\noindent --- Justice Powell, \href{https://supreme.justia.com/cases/federal/us/434/452/}{434 U.S. at 473}.

\begin{quote}
``The Compact Clause, however, is directed to joint action by more than one State. If its only purpose in the present context were to require the consent of Congress to agreements between States that would otherwise violate the Commerce Clause, it would have no independent meaning. The Clause must mean that some actions which would be permissible for individual States to undertake are not permissible for a group of States to agree to undertake.''
\end{quote}
\noindent --- Justice White, dissenting, \href{https://www.law.cornell.edu/supremecourt/text/434/452}{434 U.S. at 482}.


\end{document}
